\section{Improvement: Rewriting the Coinduction Library}
\label{sec:rewrite}

Goal was to fix bugs. But OCaml plugin got in the way.

\subsection{OCaml plugin: why it exists}

Exists to \gap{\ldots}

\note{The reason, from the library's own comments: applying tower induction needs
the goal presented as a predicate applied to the candidate, and guessing that
predicate is the hard part. The released library reified the goal into a small
syntax (hole, abstraction, conjunction) that is inf-closed by construction, which
discharges the first obligation for free. The cost is getting back a
user-readable goal: interpreting the reified syntax with Rocq's own reduction
performs unwanted simplifications \emph{and discards the user's bound-variable
names}. The plugin reconstructed the goal type in OCaml by walking the original
goal syntactically, purely to recover those names. That is all $307$ lines were
for.}

\subsection{Removing it}

\gap{\ldots}

\begin{notebox}
What replaced it: \code{pattern} already abstracts a goal over a subterm and
preserves names because it never leaves the goal's own syntax. Tower induction
becomes three lines, with the inf-closedness obligation going to proof search
instead of being free by construction:
\begin{lstlisting}
Ltac tower_induction_with c :=
  pattern (elem c);
  apply tower;
  [icauto | try clear_old_chain].
\end{lstlisting}
\end{notebox}

\subsection{Fixing bugs became much easier}

With bug fix and other error fixes, the improvement is\ldots
\gap{\ldots}

\note{The fixed bugs, for reference: \gfp{} in a premise; accumulation failing
when a premise mentioned the candidate under a function (\code{f (elem c)}); a
candidate name clashing with a binder of the goal; \code{step} not seeing through
a definition, and taking more than one step. Each is a compiling goal in
\code{paper-data/probes.v}.}

\subsection{What else did I do?}

\paragraph{Automatic monotonicity.}
In practice, 18 out of 19 of the stated monotonicity proofs in the Interaction
Trees library can be automatically discharged using \code{monauto}, with the last
being a complex proof involving multiple nontrivial steps. \code{monauto} also
discharges every one of the monotonicity sub-proofs for \code{ptower} in the
Interaction Trees library, which number \gap{\_}.

\note{\textbf{The counts have moved; current tree says 17 of 18.} There are $18$
statements of the form \code{Proper (leq ==> leq) f} in ITrees, and $17$ are
\code{Proof. monauto. Qed.} verbatim. The exception is \code{iter_mono} in
\code{extra/Dijkstra/PureITreeDijkstra.v}, whose functor is packaged in a
dependent record, so the proof must destruct it and apply a monotonicity field
stored inside before automation can start. For the blank: \code{tower induction}
is used at $22$ sites, and \code{apply_ptower'} sends its two side conditions to
\code{monauto} and \code{icauto}. ITrees contains \emph{zero} calls to
\code{icauto}, which is the sharper way to say it: the obligations never reach
the user.}

\paragraph{Automatic inf-closedness.}
\gap{\ldots}

\note{\code{icauto} is \code{auto} over a hint database. Most cases are
structural. The one that is not: for a relation class such as \code{Reflexive} or
a \code{Proper} instance, the inf-closedness is visible only after unfolding, and
the unfolding sits inside the predicate rather than at its head, so
\code{Hint Unfold} cannot reach it. A \code{Hint Extern} does that one delta
step. Claim 8 of \code{probes.v} pins this at ITrees' arity.}

\note{\textbf{Two rewrites in the git log that the outline does not mention}, and
which belong in this section next to the two above: \code{accumulate} and
\code{symmetric} were both reimplemented without the plugin. Accumulation now
reverts every hypothesis mentioning the candidate, folds them into one
conjunction, applies \code{ptower}, and restores them under their original names,
which is also what let it detect and report the one conclusion shape it cannot
support.}

\subsection{What the library cost and gained}

%% The coinduction library itself, before and after the rewrite.
\begin{table}[t]
\caption{The coinduction library before and after the rewrite, in comment-free
non-blank lines. The tactic implementation is the tactic file together with the
OCaml plugin it called. The plugin is gone, so the rewritten library has no
OCaml in it at all.}
\label{tab:library}
\centering
\small
\begin{tabular}{@{}lrrr@{}}
\toprule
Component & Released & Rewritten & Change \\
\midrule
\multicolumn{4}{@{}l}{\emph{Tactic implementation}} \\
\quad \code{tactics.v}            & 229 & 188 & $-18\%$ \\
\quad OCaml plugin                & \pluginLines{} & 0 & $-100\%$ \\
\quad total                       & \oldTacImpl{} & \newTacImpl{} & $\mathbf{-\tacImplDrop{}}$ \\
\addlinespace
\multicolumn{4}{@{}l}{\emph{Rest of the library}} \\
\quad lattice theory              & 384 & 499 & $+30\%$ \\
\quad tower                       & 134 & 150 & $+12\%$ \\
\quad companion                   & 523 & 499 & $-5\%$ \\
\quad relations                   & 124 & 148 & $+19\%$ \\
\addlinespace
Core total (excluding tests)      & \oldLibCore{} & \newLibCore{} & $+6\%$ \\
Comments                          & 370 & 455 & $+23\%$ \\
Tests                             & \oldTests{} & \newTests{} & $+290\%$ \\
\bottomrule
\end{tabular}
\end{table}


\note{This subsection is not in the outline; the table needed somewhere to live.
Delete the heading and fold \Cref{tab:library} wherever you prefer.}
